\section{Evaluation}
Before looking at throughput or scalability, the first thing I wanted to check was whether the two custom queues are actually valid and whether they behave the same way as the JDK \texttt{ConcurrentLinkedQueue}. For that reason, I first ran a set of custom correctness tests for both \texttt{MSConcurrentLinkedQueue} and \texttt{LinkedConcurrentRingQueue}. These tests compare the custom queues directly with \texttt{ConcurrentLinkedQueue} and check the basic queue API step by step. In practice this means checking operations such as \texttt{add}, \texttt{offer}, \texttt{poll}, \texttt{peek}, \texttt{contains}, \texttt{remove}, \texttt{iterator}, \texttt{spliterator}, and \texttt{toArray}, along with \texttt{null} handling and general FIFO behavior. I also used randomized differential tests with many mixed operations, so that both custom queues and the official JDK queue receive the same sequence of calls and can be compared on returned values and final contents.

I also wanted this check to be done in a setup that is as close as possible to real library use, and not only as standalone classes in my own source folder. Because of that, the two implementations were installed directly into a local OpenJDK 21 build under \texttt{java.util.concurrent} as \texttt{MSConcurrentLinkedQueue} and \texttt{LinkedConcurrentRingQueue}. The JDK build used for these tests reports itself as \texttt{OpenJDK 21.0.12-internal}. On top of the sequential and randomized API checks, I also ran multithreaded stress tests across different thread counts and small JCStress tests for simple race patterns such as two concurrent \texttt{offer} calls and \texttt{offer}/\texttt{poll} interleavings. The point of all this was to make sure that before talking about performance, the two custom queues already look correct and behave like the official JDK queue in the cases tested.

This section should define how the queue implementations will be evaluated. The evaluation should include throughput, latency, scalability across thread counts, and behavior under different producer-consumer ratios.

\subsection{Experimental Setup}
The experiments were run on my laptop, a Lenovo 20VG system with an AMD Ryzen 5 4500U with Radeon Graphics. The machine has 8 GB of DDR4 system memory, 384 KB of L1 cache, 3 MB of L2 cache, and 8 MB of L3 cache. The queue implementations were tested inside the local OpenJDK 21.0.12-internal build where the two custom queues were installed under \texttt{java.util.concurrent}.

During the experiments, the laptop was connected to power, so it was not running on battery. I also kept turbo boost disabled in order to reduce frequency changes during the measurements. This was checked with \texttt{sudo turbostat} before the runs.
\subsection{Metrics}
For the performance part, I used both some simple custom timing tests and JMH. The custom timing code is under \texttt{test/benchmarks/custom}. These tests are useful because they are easy to read and they show clearly what is being measured. In particular, I measured sequential \texttt{offer}+\texttt{poll} runs with 1,000,000 operations, and I also measured threaded \texttt{offer}/\texttt{poll} runs with 1, 4, 8, 16, and 32 threads. In the threaded tests, the code also checks things like checksum, null polls, and remaining elements, so the timing is not separated from basic correctness.

I also used JMH, the Java Microbenchmark Harness. JMH is the standard Java microbenchmark framework and it helps avoid some common benchmarking mistakes by doing warm-up, controlled measurement iterations, and repeated execution under the JVM. Here it was used mainly to measure the round-trip queue cost, where each benchmark operation does an \texttt{offer()} followed by a \texttt{poll()}, and to compare the same three implementations in a more controlled setup.

The three queues compared in these measurements are the official JDK \texttt{ConcurrentLinkedQueue}, my custom \texttt{MSConcurrentLinkedQueue}, and \texttt{LinkedConcurrentRingQueue}. For the final comparison, I report execution times in milliseconds (ms), so lower values mean better performance.

\paperfigure{figure20.png}{Initial timing results in milliseconds for the three queue implementations.}
